\documentclass[a4paper,14pt]{extarticle}

\usepackage[utf8]{inputenc}
\usepackage[russian]{babel}
\usepackage[a4paper, margin=2.5cm]{geometry}

\usepackage{amssymb}
\usepackage{amsmath}

\usepackage{setspace}

\doublespacing

\DeclareMathOperator{\M}{M}
\DeclareMathOperator{\Uniform}{U}

\begin{document}
    \begin{center}
        \textbf{Покупка машины}
    \end{center}

    M = 1\_000\_000, $S = 5\M$ -- стоимость машины, $T \in [3; 11]$ -- момент времени для продажи машины, $X \sim \Uniform(1; 11)$ -- срок службы машины, возможный доход -- 
    $\displaystyle
    G_X(x) = \begin{cases}
        1.4\M, & x \in [1; 3) \\
        0, & x \in [3; T) \\
        0.8\M, & x \in [T; 11] \\
        0, & x \notin [1; 11]
    \end{cases} = \\
    1.4\M \cdot I(x \in [1; 3)) + 0.8\M \cdot I(x \in [T; 11])
    $

    Случайная величина $\displaystyle \xi = \frac{S - G_X(x)}{\min(X, T)}$

    Плотность вероятностей $\displaystyle X: f_X(x) = \frac{I(x \in [1; 11])}{10}$

    Функция распределения $\displaystyle X: F_X(x) = \frac{x - 1}{10} \cdot I(x \in [1; 11]) + I(x > 11)$\\
    $\displaystyle
    \text{При } x \in [1; 3): \\
    G_X(x) = 1.4\M \\
    \xi_1 = \frac{5\M - 1.4\M}{x} \cdot I(x \in [1; 3)) = \frac{3.6\M}{x} \cdot I(x \in [1; 3)) \\
    \text{При } x \in [3; T): \\
    G_X(x) = 0 \\
    \xi_2 = \frac{5\M - 0}{x} \cdot I(x \in [3; T)) = \frac{5\M}{x} \cdot I(x \in [3; T)) \\
    \text{При } x \in [T; 11]: \\
    G_X(x) = 0.8\M \\
    \xi_3 = \frac{5\M - 0.8\M}{T} \cdot I(x \in [T; 11]) = \frac{4.2\M}{T} \cdot I(x \in [T; 11]) \\
    $

    Таким образом математическое ожидание равно:\\
    $\displaystyle
    E(\xi) = \sum_{i = 1}^{3} \int_{\mathbf{R}} \xi_i \cdot f_X(x) \,dx = \int_{1}^{3} \frac{3.6\M}{x} \cdot \frac{1}{10} \,dx + \int_{3}^{T} \frac{5\M}{x} \cdot \frac{1}{10} \,dx + \int_{T}^{11} \frac{4.2\M}{T} \cdot \frac{1}{10} \,dx = 
    0.36\M \cdot \ln(x)|_{1}^{3} + 0.5\M \cdot \ln(x)|_{3}^{T} + \frac{0.42\M}{T} \cdot x|_{T}^{11} = 
    0.36\M \cdot \ln3 + 0.5\M \cdot \ln T - 0.5\M \cdot \ln3 + \frac{0.42\M}{T}(11 - T) = -0,14\M \cdot \ln3 + 0.5\M \cdot \ln T + \frac{4.62\M - 0.42\M \cdot T}{T}
    $

    Для поиска минимального $T$ нужно найти точку минимума для $E(\xi)(T)$:\\
    $\displaystyle
    \frac{d}{dT} E(\xi)(T) = \frac{0.5}{T} - \frac{4.62\M}{T^2} = 0 \\
    T_{\min} = 9.24
    $

\end{document}